% ============================================================
\section{Esercizi per l'esame}
\label{sec:esercizi-esame}
% ============================================================

\begin{softnote}
	Questa sezione raccoglie esercizi tipici d'esame con soluzione completa:
	inserimento in ABR, analisi di complessità su grafi sparsi, valutazione di polinomi
	e ricerca binaria ricorsiva con modelli di costo diversi per il passaggio dei parametri.
\end{softnote}

% ------------------------------------------------------------
\subsection{Esercizio 1 --- Inserimento in ABR con chiavi duplicate}
\label{sub:es1-abr-insert}
% ------------------------------------------------------------

Si considerino gli alberi binari di ricerca e la seguente procedura \textsc{Insert}.

\begin{algorithm}[H]
	\caption{\textsc{Insert}$(t,\, z)$ --- versione standard (riferimento)}
	\KwIn{Un ABR $t$ e un nodo $z \notin t$}
	$y \leftarrow \textsc{null}$\;
	$x \leftarrow t.\mathit{root}$\;
	\While{$x \neq \textsc{null}$}{
		$y \leftarrow x$\;
		$x \leftarrow (z.\mathit{key} < x.\mathit{key})$
		\textbf{ ? } $x.\mathit{left}$ \textbf{ : } $x.\mathit{right}$\;
	}
	$z.p \leftarrow y$\;
	\eIf{$y = \textsc{null}$}{
		$t.\mathit{root} \leftarrow z$\;
	}{
		\eIf{$z.\mathit{key} < y.\mathit{key}$}{
			$y.\mathit{left} \leftarrow z$\;
		}{
			$y.\mathit{right} \leftarrow z$\;
		}
	}
\end{algorithm}

\paragraph{Domanda.}
Si determinino le prestazioni asintotiche della procedura \textsc{Insert} quando si
inseriscono $n$ chiavi \emph{identiche} in un ABR inizialmente vuoto, per ciascuna delle
tre varianti seguenti.

\paragraph{Variante A (standard).}
Come nella procedura sopra: chiavi uguali vanno sempre a destra.

\paragraph{Variante B (flag booleano).}
In ogni nodo $x$ è presente un flag booleano $b[x]$. La riga di discesa diventa:
\begin{enumerate}
	\item $b[x] \leftarrow \lnot b[x]$;
	\item se $z.\mathit{key} = x.\mathit{key}$, allora
	$x \leftarrow b[x]\ \textbf{?}\ x.\mathit{left}\ \textbf{:}\ x.\mathit{right}$;
	\item altrimenti $x \leftarrow (z.\mathit{key} < x.\mathit{key})\ \textbf{?}\ x.\mathit{left}\ \textbf{:}\ x.\mathit{right}$.
\end{enumerate}
Le righe di inserimento (11--12) vanno modificate in modo coerente: quando
$z.\mathit{key} = y.\mathit{key}$, il nuovo nodo $z$ deve essere inserito nel figlio
scelto nell'ultima iterazione del ciclo, \emph{senza} modificare nuovamente $b[y]$.

\paragraph{Variante C (lista di duplicati).}
Ogni nodo mantiene una lista delle chiavi uguali alla propria. Se $z.\mathit{key} = x.\mathit{key}$,
$z$ non diventa un nuovo nodo dell'albero ma viene aggiunto alla lista di $x$.

\subsubsection*{Soluzione}

\paragraph{1. Procedura \textsc{Insert} standard.}
Il test alla riga di discesa è $z.\mathit{key} < x.\mathit{key}$. Se le due chiavi sono
uguali, il test è falso e la discesa prosegue sempre verso destra. Ogni nuova chiave viene
inserita come figlio destro del nodo più a destra: l'albero degenera in una lista.

L'$i$-esimo inserimento percorre una catena di lunghezza $i-1$, quindi costa $\Theta(i)$.
Il costo totale per inserire $n$ chiavi identiche è
\[
\sum_{i=1}^{n} \Theta(i) = \Theta(n^2).
\]
L'albero finale ha altezza $\Theta(n)$.

\paragraph{2. Procedura \textsc{Insert} con flag booleano.}
Quando si incontra una chiave uguale a quella del nodo corrente $x$, il flag $b[x]$ viene
invertito e la discesa alterna tra sottoalbero sinistro e destro. Questo avviene in ogni
nodo visitato con chiave uguale, quindi le chiavi duplicate si distribuiscono in modo
bilanciato tra i due sottoalberi.

Per induzione, l'albero costruito inserendo $n$ chiavi identiche ha altezza $\Theta(\log n)$.
Il costo dell'$i$-esimo inserimento è $\Theta(\log i)$. Il costo totale è
\[
\sum_{i=1}^{n} \Theta(\log i) = \Theta(n \log n).
\]
L'albero finale ha altezza $\Theta(\log n)$.

\paragraph{3. Procedura \textsc{Insert} con lista di duplicati.}
La prima chiave diventa la radice. Tutte le successive sono uguali alla chiave della radice:
non si prosegue la discesa e ogni inserimento aggiorna solo la lista associata alla radice.

Se l'inserimento nella lista avviene in testa, ogni inserimento nella lista costa $\Theta(1)$
e il costo totale è $\Theta(n)$. L'albero finale contiene un solo nodo (la radice) con una
lista di $n-1$ elementi.

\begin{osservazione}
	Se invece si inserisse in coda alla lista senza un puntatore all'ultimo elemento, il
	costo totale potrebbe diventare $\Theta(n^2)$.
\end{osservazione}

% ------------------------------------------------------------
\subsection{Esercizio 2 --- \texttt{edges()} in \texttt{SparseWeightedGraph}}
\label{sub:es2-sparse-graph}
% ------------------------------------------------------------

Si consideri il seguente metodo della classe \texttt{SparseWeightedGraph}:

\begin{lstlisting}[language=Java, caption={Versione originale di \texttt{edges()}}]
public List<E> edges() {
  List<E> edges = new ArrayList<>(size());
  for (MainEntry nodeE : nodes) {
    for (AdjacentEntry pair : nodeE.adjacent) {
      if (edges.contains(pair.edge.getName())) {
        continue;
      }
      edges.add(pair.edge);
    }
  }
  return edges;
}
\end{lstlisting}

La classe rappresenta il grafo con liste di adiacenza. I nodi hanno nomi di tipo
\texttt{String} che sono identificatori nel grafo. Il metodo \texttt{List.contains} scorre
tutta la lista per la ricerca.

\paragraph{Domanda.}
\begin{enumerate}
	\item Determinare la complessità in tempo rispetto al numero di nodi $n$ e di archi $m$.
	\item Modificare la procedura in modo che la complessità in tempo sia $O(n+m)$.
\end{enumerate}

\subsubsection*{Soluzione}

Il ciclo esterno (riga 3) esegue $n$ iterazioni. Il ciclo interno (riga 4), sommando gli
adiacenti di tutti i nodi, visita in totale $2m$ coppie (ogni arco non orientato compare
in due liste). Senza la riga 5 il costo sarebbe $O(n+m)$.

L'istruzione \texttt{edges.contains(...)} obbliga a verificare che l'arco non sia già stato
inserito. Se alla fine della lista ci sono $k$ archi, ogni \texttt{contains} costa $\Theta(k)$.
Nel caso peggiore si ottiene
\[
\sum_{i=1}^{m} i + \sum_{i=m+1}^{2m} m = O(m^2),
\]
quindi il costo finale è $O(n + m^2)$.

\paragraph{Modifica in $O(n+m)$.}
Un arco non orientato tra $A$ e $B$ compare in entrambe le liste di adiacenza. Poiché i nomi
sono confrontabili (\texttt{String}), si inserisce l'arco solo quando il nome del primo
estremo precede (lessicograficamente) l'altro:

\begin{lstlisting}[language=Java, caption={Versione ottimizzata di \texttt{edges()}}]
public List<E> edges() {
  List<E> edges = new ArrayList<>(size());
  for (MainEntry nodeE : nodes) {
    for (AdjacentEntry pair : nodeE.adjacent) {
      if (pair.node.getName().compareTo(nodeE.node.getName()) <= 0) {
        continue;
      }
      edges.add(pair.edge);
    }
  }
  return edges;
}
\end{lstlisting}

Ogni arco viene considerato esattamente una volta: costo $O(n+m)$.

% ------------------------------------------------------------
\subsection{Esercizio 3 --- Valutazione di polinomi (Horner vs sommatoria)}
\label{sub:es3-horner}
% ------------------------------------------------------------

Si vuole determinare il valore del polinomio
\[
P(x) = \sum_{k=0}^{n} a_k x^k
\]
usando la regola di Horner:
\[
P(x) = a_0 + x\bigl(a_1 + x(a_2 + \cdots + x(a_{n-1} + x a_n)\cdots)\bigr).
\]
I coefficienti $a_0,\ldots,a_n$ e un valore di $x$ sono forniti come input.

\paragraph{Domanda.}
\begin{enumerate}
	\item Scrivere un metodo Java che calcoli $P(x)$ con la regola di Horner e determinare
	il tempo asintotico in $\Theta$.
	\item Scrivere un metodo Java che calcoli $P(x)$ dalla sommatoria, calcolando ciascun
	termine $a_i x^i$ separatamente. Per ogni $i$, la potenza $x^i$ deve essere calcolata
	da zero, senza riusare valori intermedi. Determinare il tempo asintotico in $\Theta$.
\end{enumerate}

\subsubsection*{Soluzione}

\paragraph{a) Regola di Horner.}
Il calcolo parte dal coefficiente di grado massimo $a_n$ e procede verso $a_0$.

\begin{lstlisting}[language=Java, caption={Valutazione con la regola di Horner}]
public static double horner(double[] a, double x) {
    double p = a[a.length - 1];
    for (int i = a.length - 2; i >= 0; i--) {
        p = a[i] + x * p;
    }
    return p;
}
\end{lstlisting}

Se il vettore \texttt{a} contiene $n+1$ coefficienti, il ciclo viene eseguito $n$ volte con
un numero costante di operazioni aritmetiche per iterazione. Il tempo di computazione è
$\Theta(n)$.

\paragraph{b) Sommatoria diretta senza riuso delle potenze.}

\begin{lstlisting}[language=Java, caption={Valutazione diretta senza riuso}]
public static double diretto(double[] a, double x) {
    double p = 0.0;
    for (int i = 0; i < a.length; i++) {
        double potenza = 1.0;
        for (int j = 0; j < i; j++) {
            potenza = potenza * x;
        }
        p = p + a[i] * potenza;
    }
    return p;
}
\end{lstlisting}

Per calcolare $a_i x^i$, il ciclo interno esegue $i$ moltiplicazioni. Il costo totale è
\[
\sum_{i=0}^{n} \Theta(i) = \Theta\!\left(\sum_{i=0}^{n} i\right) = \Theta\!\left(\frac{n(n+1)}{2}\right) = \Theta(n^2).
\]

\paragraph{Confronto.}
\begin{itemize}
	\item Horner: $\Theta(n)$.
	\item Metodo diretto senza riuso: $\Theta(n^2)$.
\end{itemize}
La regola di Horner è asintoticamente più efficiente.

% ------------------------------------------------------------
\subsection{Esercizio 4 --- Ricerca binaria ricorsiva e costo del passaggio parametri}
\label{sub:es4-binsearch-params}
% ------------------------------------------------------------

Normalmente si assume che il passaggio di un array di $n$ elementi a una procedura costi
$\Theta(1)$ (riferimento). Si considerino tre ipotesi alternative:

\begin{enumerate}
	\item L'array viene passato tramite \textbf{riferimento}: costo $\Theta(1)$.
	\item L'array viene passato con \textbf{copia completa}: costo $\Theta(n)$.
	\item Viene passato solo il \textbf{sottoarray} $A[p..q]$ corrente: costo $\Theta(q-p+1)$.
\end{enumerate}

Si consideri la ricerca binaria ricorsiva su un array ordinato di dimensione $n$. Al livello
$i$ della ricorsione il sottoproblema ha dimensione circa $n/2^i$. Escludendo il passaggio,
il lavoro a ogni chiamata è $\Theta(1)$.

\paragraph{Domanda.}
Per ciascuna ipotesi, determinare l'equazione di ricorrenza del tempo nel caso peggiore e
una soluzione asintotica stretta in $\Theta$.

\subsubsection*{Soluzione}

A ogni chiamata la ricerca prosegue su una sola metà: la dimensione del sottoproblema si dimezza.

\paragraph{a) Passaggio per riferimento.}
Il costo di passaggio è $\Theta(1)$. Ricorrenza:
\[
T(n) = T(n/2) + \Theta(1), \qquad T(1) = \Theta(1).
\]
Profondità $\Theta(\log n)$ con costo $\Theta(1)$ per livello. Quindi $T(n) = \Theta(\log n)$.

\paragraph{b) Copia completa dell'array a ogni chiamata.}
A ogni livello si copia l'intero array iniziale di dimensione $n$, anche se il sottoproblema
logico ha dimensione $n/2^i$. Indicando con $T_i(n)$ il tempo residuo dal livello $i$:
\[
T_i(n) = T_{i+1}(n) + \Theta(n),
\]
con profondità $\Theta(\log n)$. Ci sono $\Theta(\log n)$ livelli, ciascuno con costo
$\Theta(n)$. Quindi $T(n) = \Theta(n \log n)$.

\paragraph{c) Passaggio del solo sottoarray $A[p..q]$.}
Se il sottoproblema corrente ha dimensione $n$, il costo di passaggio del sottoarray è
$\Theta(n)$. Ricorrenza:
\[
T(n) = T(n/2) + \Theta(n), \qquad T(1) = \Theta(1).
\]
Sviluppando:
\[
T(n) = \Theta(n) + \Theta(n/2) + \Theta(n/4) + \cdots + \Theta(1).
\]
La somma $n + n/2 + n/4 + \cdots + 1$ è una serie geometrica di ragione $1/2$:
\[
n \leq n + \frac{n}{2} + \frac{n}{4} + \cdots + 1 < 2n,
\]
quindi la somma è $\Theta(n)$ e $T(n) = \Theta(n)$.

\paragraph{Confronto finale.}
\begin{itemize}
	\item Riferimento: $T(n) = \Theta(\log n)$.
	\item Copia completa a ogni chiamata: $T(n) = \Theta(n \log n)$.
	\item Solo sottoarray corrente: $T(n) = \Theta(n)$.
\end{itemize}
Il modello di costo per il passaggio dei parametri può modificare in modo significativo il
tempo asintotico di una procedura ricorsiva apparentemente identica.
